% The action of the transformation $t-\lambda_i$ on $\gennullspace{t-\lambda_i}$
% is especially easy to understand.
% Observe that any transformation~$t$ is nilpotent on $\gennullspace{t}$, 
% because if $\vec{v}\in\gennullspace{t}$ then 
% by definition $t^n(\vec{v})=\zero$. 
% Thus $t-\lambda_i$ is nilpotent on 
% $\gennullspace{t-\lambda_i}$.

% % Thus, the generalized null space $\gennullspace{t-\lambda_i}$ is a part of
% % the space on which the action of $t-\lambda_i$ is easy to understand.

% We shall take three steps to prove this section's major result. 
% The next result is the first.
% % It establishes that if $j=neq i$ then \( t-\lambda_j \) leaves
% % \( t-\lambda_i \)'s part unchanged.

% \begin{lemma} \label{le:tInvIfftMinLambdaInv}
% A subspace is \( t \) invariant if and only if 
% it is \( t-\lambda \) invariant for all scalars \( \lambda \).
% In particular, 
% if \( \lambda_i \) is an eigenvalue of  a linear transformation
% \( t \) then for any other eigenvalue $\lambda_j$
% the spaces \( \gennullspace{t-\lambda_i} \) 
% and \( \genrangespace{t-\lambda_i} \)
% are \( t-\lambda_j \) invariant.
% \end{lemma}

% \begin{proof}
% For the first sentence we check the two implications separately.
% The `if' half is easy: if the subspace is $t-\lambda$ invariant for 
% all scalars $\lambda$ then using $\lambda=0$ shows that it is $t$ invariant.
% For `only if' suppose that the subspace is $t$ invariant,
% so that if $\vec{m}\in M$ then $t(\vec{m})\in M$, and let $\lambda$ be 
% a scalar.
% The subspace $M$ is closed under linear combinations and so if 
% $t(\vec{m})\in M$ then $t(\vec{m})-\lambda\vec{m}\in M$.
% Thus if $\vec{m}\in M$ then $(t-\lambda)\,(\vec{m})\in M$.

% The lemma's second sentence follows from its first.
% The
% two spaces are $t-\lambda_i$~invariant so they are \( t \)~invariant.
% Apply the first sentence again to
% conclude that they are also \( t-\lambda_j \) invariant.
% \end{proof}

% The second step of the three that we will take to prove this
% section's major result makes use of an additional property of 
% \( \gennullspace{t-\lambda_i} \) and
% \( \genrangespace{t-\lambda_i} \), that they are complementary.
% Recall that if a space is the direct sum of two others 
% \( V=\mathscr{N}\directsum \mathscr{R} \) 
% then any vector \( \vec{v} \) in the space breaks into
% two parts \( \vec{v}=\vec{n}+\vec{r} \) where \( \vec{n}\in \mathscr{N} \) and
% \( \vec{r}\in \mathscr{R} \), and recall also 
% that if \( B_{\mathscr{N}} \) and \( B_{\mathscr{R}} \) are bases for
% \( \mathscr{N} \) and \( \mathscr{R} \) then the concatenation
% \( \cat{B_{\mathscr{N}}}{B_{\mathscr{R}}} \) is linearly independent.
% The next result says that for any subspaces
% \( \mathscr{N} \) and \( \mathscr{R} \) that are complementary 
% as well as \( t \)~invariant,
% the action
% of \( t \) on \( \vec{v} \) breaks into the actions of
% \( t \) on \( \vec{n} \) and on \( \vec{r} \).

% \begin{lemma} \label{le:InvCompSubspSplitTrans}
% Let \( \map{t}{V}{V} \) be a transformation and let \( \mathscr{N} \) and 
% \( \mathscr{R} \) be
% \( t \) invariant complementary subspaces of \( V \).
% Then we can represent \( t \) by a matrix with
% blocks of square submatrices $T_1$ and $T_2$
% \begin{equation*} \renewcommand{\arraystretch}{1.2}
%   \begin{pmat}{c|c}
%       T_1   &Z_2  \\  \cline{1-2}
%       Z_1 &T_2
%    \end{pmat}
%    \begin{array}{@{}l}
%      \} \text{\ $\dim(\mathscr{N})$-many rows}  \\
%      \} \text{\ $\dim(\mathscr{R})$-many rows}
%    \end{array}
% \end{equation*}
% where \( Z_1 \) and \( Z_2 \) are blocks of zeroes.
% \end{lemma}

% \begin{proof}
% Since the two subspaces are complementary, the concatenation of a basis
% for \( \mathscr{N} \) with a basis for \( \mathscr{R} \) makes a basis
% \( B=\sequence{\vec{\nu}_1,\dots,\vec{\nu}_p,
%         \vec{\mu}_1,\ldots,\vec{\mu}_q}  \)
% for \( V \).
% We shall show that the matrix
% \begin{equation*}
%   \rep{t}{B,B}=
%   \begin{pmat}{c@{\hspace*{1em}}c@{\hspace*{1em}}c}
%      \vdots                   &        &\vdots     \\
%      \rep{t(\vec{\nu}_1)}{B}  &\cdots  &\rep{t(\vec{\mu}_q)}{B}  \\
%      \vdots                   &        &\vdots     \\
%   \end{pmat}
% \end{equation*}
% has the desired form.

% Any vector \( \vec{v}\in V \) is a member of \( \mathscr{N} \) 
% if and only if when it is represented with respect to \( B \)
% the final \( q \)
% coefficients are zero.
% As \( \mathscr{N} \) is \( t \)~invariant, each of the vectors
% \( \rep{t(\vec{\nu}_1)}{B} \),
% \ldots, \( \rep{t(\vec{\nu}_p)}{B} \) has this form.
% Hence the lower left of \( \rep{t}{B,B} \) is all zeroes.
% The argument for the upper right is similar.
% \end{proof}

% To see that we have decomposed \( t \) into its action on the parts, 
% let $B_{\mathscr{N}}=\sequence{\vec{\nu}_1,\dots,\vec{\nu}_p}$ and 
% $B_{\mathscr{R}}=\sequence{\vec{\mu}_1,\ldots,\vec{\mu}_q}$.
% The restrictions of \( t \) to the subspaces \( \mathscr{N} \) 
% and~\( \mathscr{R} \) 
% are represented
% with respect to the bases $B_{\mathscr{N}},B_{\mathscr{N}}$
% and $B_{\mathscr{R}},B_{\mathscr{R}}$
% by the matrices \( T_1 \) and \( T_2 \).
% So with subspaces that are invariant and complementary 
% we can split the problem of examining
% a linear transformation into two lower-dimensional subproblems.
% The next result illustrates this decomposition into blocks.

% \begin{lemma} \label{le:DetIsProdOfSubDets}
% If $T$ is a matrix with square submatrices $T_1$ and $T_2$
% \begin{equation*} \renewcommand{\arraystretch}{1.2}
%   T=
%   \begin{pmat}{c|c}
%       T_1   &Z_2  \\  \cline{1-2}
%       Z_1   &T_2
%    \end{pmat}
% \end{equation*}
% where the \( Z \)'s are blocks of zeroes,
% then \( \deter{T}=\deter{T_1}\cdot\deter{T_2} \).
% \end{lemma}

% \begin{proof}
% Suppose that \( T \) is \( \nbyn{n} \),
% that \( T_1 \) is \( \nbyn{p} \),
% and that \( T_2 \) is \( \nbyn{q} \).
% In the permutation formula for the determinant
% \begin{equation*}
%   \deter{T}=
%   \sum_{\text{\scriptsize permutations\ }\phi}
%           t_{1,\phi(1)}t_{2,\phi(2)}\cdots t_{n,\phi(n)}\sgn(\phi)
% \end{equation*}
% each term comes from a rearrangement of the column numbers
% \( 1,\dots,n \) into a new order \( \phi(1),\dots,\phi(n) \).
% The upper right block $Z_2$ is all zeroes, so if a
% \( \phi \) has at least one of \( p+1,\dots,n \) among its first
% \( p \) column numbers \( \phi(1),\dots,\phi(p) \) then the term arising
% from \( \phi \) does not contribute to the sum because it is zero,
% e.g., if \( \phi(1)=n \) then
% \( t_{1,\phi(1)}t_{2,\phi(2)}\dots t_{n,\phi(n)}
%    =0\cdot t_{2,\phi(2)}\dots t_{n,\phi(n)}=0 \).

% So the above formula reduces to a sum over all permutations with 
% two halves:~any contributing $\phi$ is the composition of a $\phi_1$ that
% rearranges only \( 1,\dots,p \) 
% and a $\phi_2$ that rearranges only \( p+1,\dots,p+q \).
% Now, the distributive law 
% and the fact that the signum of a composition is the product
% of the signums gives that this
% \begin{multline*}
%    \deter{T_1}\cdot\deter{T_2}=
%    \bigg(\sum_{\begin{subarray}{c}
%                 \text{\scriptsize perms\ }\phi_1 \\
%                 \text{\scriptsize of\ } 1,\dots,p
%                \end{subarray}}
%        \!\!\! t_{1,\phi_1(1)}\cdots t_{p,\phi_1(p)}\sgn(\phi_1) \bigg)  \\
%    \cdot
%    \bigg(\sum_{\begin{subarray}{c}
%                 \text{\scriptsize perms\ }\phi_2 \\
%                 \text{\scriptsize of\ } p+1,\dots,p+q
%                \end{subarray}}
%        \!\!\! t_{p+1,\phi_2(p+1)}\cdots t_{p+q,\phi_2(p+q)}\sgn(\phi_2) 
%         \bigg)
% \end{multline*}
% equals
%   $\deter{T}=
%   \sum_{\text{\scriptsize contributing\ }\phi}
%           t_{1,\phi(1)}t_{2,\phi(2)}\cdots t_{n,\phi(n)}\sgn(\phi)$.
% \end{proof}

% \begin{example}
% \begin{equation*}
%     \begin{vmat}[r]
%        2  &0  &0  &0  \\
%        1  &2  &0  &0  \\
%        0  &0  &3  &0  \\
%        0  &0  &0  &3
%     \end{vmat}
%    =\begin{vmat}[r]
%        2  &0  \\
%        1  &2
%     \end{vmat}
%     \cdot
%     \begin{vmat}[r]
%        3  &0  \\
%        0  &3
%     \end{vmat}
%    =36
% \end{equation*}
% \end{example}

% From \nearbylemma{le:DetIsProdOfSubDets} we conclude that
% if two subspaces 
% are complementary and \( t \)~invariant then
% \( t \) is one-to-one if and only if its 
% restriction %\appendrefs{restrictions}
% to each subspace is nonsingular.

% Now for the promised third, and final, step to the main result.

% \begin{lemma}
% If a linear transformation \( \map{t}{V}{V} \) has the 
% characteristic polynomial
% \( (x-\lambda_1)^{p_1}\dots(x-\lambda_k)^{p_k} \) then
% (1)~\( V=\gennullspace{t-\lambda_1}\directsum\cdots
%              \directsum\gennullspace{t-\lambda_k} \) 
% and
% (2)~\( \dim(\gennullspace{t-\lambda_i})=p_i  \).
% \end{lemma}

% \begin{proof}
% This argument consists of proving
% two preliminary claims, followed by proofs of clauses~(1) and~(2).

% The first claim is that
% \( \gennullspace{t-\lambda_i}\intersection\gennullspace{t-\lambda_j}=\set{\zero} \)
% when \( i\neq j \).
% By \nearbylemma{le:tInvIfftMinLambdaInv}
% both \( \gennullspace{t-\lambda_i} \) and 
% \( \gennullspace{t-\lambda_j} \) are \( t \)~invariant.
% The intersection of \( t \) invariant subspaces is \( t \)
% invariant and so the restriction of \( t \) to
% \( \gennullspace{t-\lambda_i}\intersection\gennullspace{t-\lambda_j} \)
% is a linear transformation.
% Now,
% $t-\lambda_i$ is nilpotent on \( \gennullspace{t-\lambda_i} \)
% and  
% $t-\lambda_j$ is nilpotent on \( \gennullspace{t-\lambda_j} \),
% so both \( t-\lambda_i \) and \( t-\lambda_j \) are 
% nilpotent on the intersection.
% Therefore by \nearbylemma{le:NilIffOnlyEigenZero} and the observation
% following it, if \( t \) has any eigenvalues on the intersection 
% then the ``only'' eigenvalue is both
% \( \lambda_i \) and \( \lambda_j \).
% This cannot be, so the restriction has no eigenvalues:
% \( \gennullspace{t-\lambda_i}\intersection\gennullspace{t-\lambda_j} \)
% is the trivial space
% (\nearbylemma{le:MapNonTrivSpHasEigen} shows that 
% the only transformation that is without any eigenvalues is 
% the transformation on the trivial space).

% The second claim is that  when \( i\neq j \) then
% \( \gennullspace{t-\lambda_i} \subseteq \genrangespace{t-\lambda_j} \).
% To verify it we will show that \( t-\lambda_j \) is one-to-one on
% \( \gennullspace{t-\lambda_i} \).
% This will imply, 
% since \( \gennullspace{t-\lambda_i} \) is \( t-\lambda_j \) invariant
% by \nearbylemma{le:tInvIfftMinLambdaInv},
% that the map \( t-\lambda_j \) is an automorphism of the subspace
% \( \gennullspace{t-\lambda_i} \),
% and therefore that 
% \( \gennullspace{t-\lambda_i} \)  is a subset of each
% \( \rangespace{t-\lambda_j} \), \( \rangespace{(t-\lambda_j)^2} \), etc.
% To verify that the map is one-to-one,
% suppose that \( \vec{v}\in\gennullspace{t-\lambda_i} \) is in the 
% null space of \( t-\lambda_j \), aiming to show that \( \vec{v}=\zero \).
% Consider the map 
% \( [(t-\lambda_i)-(t-\lambda_j)]^n \).
% On the one hand, the only
% vector that \( (t-\lambda_i)-(t-\lambda_j)=\lambda_i-\lambda_j \) maps to 
% zero is the zero vector itself.
% On the other hand, as in the proof of 
% \nearbylemma{le:PolyMapsFactor}
% the binomial expansion gives this.
% \begin{equation*}
%   (t-\lambda_i)^n(\vec{v})
%     +\binom{n}{1}(t-\lambda_i)^{n-1}(t-\lambda_j)^1(\vec{v})
%     +\binom{n}{2}(t-\lambda_i)^{n-2}(t-\lambda_j)^2(\vec{v})
%     +\cdots 
% \end{equation*}
% The first term is the zero vector 
% because \( \vec{v}\in\gennullspace{t-\lambda_i} \).
% The remaining terms are the zero vector because
% \( \vec{v} \) is in the null space of \( t-\lambda_j \).
% Therefore \( \vec{v}=\zero \). 

% With those two preliminary claims done 
% we can prove clause~(1), that the space is the direct sum of the
% generalized null spaces.
% By Corollary~III.\ref{GenRngNullDirSumToSp} the space is the direct sum
% \( V=\gennullspace{t-\lambda_1}\directsum\genrangespace{t-\lambda_1} \).
% By the second claim
% \( \gennullspace{t-\lambda_2}\subseteq\genrangespace{t-\lambda_1} \)
% and so we can get a basis for \( \genrangespace{t-\lambda_1} \) by
% starting with a basis for \( \gennullspace{t-\lambda_2} \) and adding
% extra basis elements taken from
% \( \genrangespace{t-\lambda_1}\intersection\genrangespace{t-\lambda_2} \). 
% Thus \( V=\gennullspace{t-\lambda_1}\directsum\gennullspace{t-\lambda_2}
%           \directsum 
%        (\genrangespace{t-\lambda_1}\intersection\genrangespace{t-\lambda_2}) \).
% Continuing in this way we get this.
% \begin{equation*}
%   V=\gennullspace{t-\lambda_1}\directsum\cdots\directsum\gennullspace{t-\lambda_k}
%           \directsum 
%        (\genrangespace{t-\lambda_1}\intersection\cdots\intersection\genrangespace{t-\lambda_k})
% \end{equation*}
% The first claim above shows that the final space is trivial. 

% We finish by verifying clause~(2).
% Decompose \( V \) as 
% \( \gennullspace{t-\lambda_i}\directsum\genrangespace{t-\lambda_i} \)
% %with basis $B=\cat{B_{\mathscr{N}}}{B_{\mathscr{R}}}$
% and apply \nearbylemma{le:InvCompSubspSplitTrans}.
% \begin{equation*}   \renewcommand{\arraystretch}{1.2}
%   T=
% %  \rep{t}{B,B}=
%   \begin{pmat}{c|c}
%       T_1   &Z_2  \\  \cline{1-2}
%       Z_1   &T_2
%    \end{pmat}
%    \begin{array}{@{}l}
%      \} \text{\ $\dim(\,\gennullspace{t-\lambda_i}\,)$-many rows}  \\
%      \} \text{\ $\dim(\,\genrangespace{t-\lambda_i}\,)$-many rows}
%    \end{array}
% \end{equation*}
% \nearbylemma{le:DetIsProdOfSubDets} says that
% \( \deter{T-xI}=\deter{T_1-xI}\cdot\deter{T_2-xI} \).
% By the uniqueness clause of the Fundamental Theorem of Algebra, 
% Theorem~I.\ref{th:FundThmAlg},
% the determinants of the blocks have the same factors as the
% characteristic polynomial
% \( \deter{T_1-xI}=(x-\lambda_1)^{q_1}\cdots(x-\lambda_z)^{q_k} \)
% and
% \( \deter{T_2-xI}=(x-\lambda_1)^{r_1}\cdots(x-\lambda_z)^{r_k} \),
% where
% \( q_1+r_1=p_1 \), \dots, \( q_k+r_k=p_k \).
% We will finish by establishing that (i)~$q_j=0$ for all $j\neq i$,
% and (ii)~$q_i=p_i$.
% Together these prove clause~(2) because they show that the degree of 
% the polynomial $\deter{T_1-xI}$ is $q_i$ and
% the degree of that polynomial equals 
% the dimension of the
% generalized null space $\gennullspace{t-\lambda_i}$. 

% For (i),
% because the restriction of \( t-\lambda_i \) to \( \gennullspace{t-\lambda_i} \)
% is nilpotent on that space,
% $t$'s only eigenvalue on that space is \( \lambda_i \),
% by \nearbylemma{cor:tMinLambdaNilpotent}.
% So $q_j=0$ for $j\neq i$.

% For (ii),
% consider the restriction of \( t \) to \( \genrangespace{t-\lambda_i} \).
% By Lemma~III.\ref{lem:RestONeToOne}, the map
% \( t-\lambda_i \) is one-to-one on
% \( \genrangespace{t-\lambda_i} \) and so \( \lambda_i \) is not an
% eigenvalue of \( t \) on that subspace.
% Therefore \( x-\lambda_i \) is not a factor of \( \deter{T_2-xI} \),
% so $r_i=0$, and so \( q_i=p_i \).
% \end{proof}

% Recall the goal of this chapter, to give a canonical form for matrix similarity.
% That result is next.
% It 
% translates the above steps into matrix terms.

% \begin{theorem}
% \index{Jordan form!represents similarity classes}\index{similar!canonical form}
% \index{canonical form!for similarity}\index{transformation!Jordan form for}
% Any square matrix is similar to one in \definend{Jordan form}
% \begin{equation*}
%   \begin{mat}
%     J_{\lambda_1}  &            &\text{\textit{--zeroes--}}                 \\
%                &J_{\lambda_2}                                              \\
%                &     &\ddots                                     \\
%                 &     &                           &J_{\lambda_{k-1}} &     \\
%                &     &\text{\textit{--zeroes--}} &            &J_{\lambda_{k}}
%   \end{mat}
% \end{equation*}
% where each \( J_{\lambda} \) is the Jordan block associated with an
% eigenvalue $\lambda$ of the original matrix (that is, each \( J_{\lambda} \)
% is all zeroes except for
% \( \lambda \)'s down the diagonal and some subdiagonal ones).
% \end{theorem}

% \begin{proof}
% Given an \( \nbyn{n} \) matrix \( T \), consider the linear map
% \( \map{t}{\C^n}{\C^n} \) that it represents
% with respect to the standard bases.
% Use the prior lemma to write
% \( \C^n=\gennullspace{t-\lambda_1}\directsum\cdots
%         \directsum\gennullspace{t-\lambda_k} \)
% where \( \lambda_1 \), \ldots, \(\lambda_k \) are the eigenvalues of \( t \).
% Because each \( \gennullspace{t-\lambda_i} \)  is \( t \) invariant,
% \nearbylemma{le:InvCompSubspSplitTrans} and the prior lemma show
% that \( t \) is represented by a matrix that is all zeroes except for square
% blocks along the diagonal.
% To make those blocks into Jordan blocks, pick each \( B_{\lambda_i} \)
% to be a string basis for the action of \( t-\lambda_i \) on
% \( \gennullspace{t-\lambda_i} \). 
% \end{proof}

% \begin{corollary}
% Every square matrix is similar to the sum of a diagonal matrix and a nilpotent
% matrix.
% \end{corollary}

\begin{remark}
严格地说，要作为矩阵相似\index{representative!of similarity classes}
的标准形，Jordan 形式必须是唯一的。
也就是说，对任意方阵，应当有且仅有一个与它相似的矩阵处于
Jordan 形式。
我们可以让定理中的形式唯一化：
将各个 Jordan 块按特征值的某个指定顺序排列，
然后再将副对角线上的 1 构成的块从最长到最短排列。
在下面的例子中，我们不讨论这一点。
\end{remark}

% As in the nilpotent cases earlier in this section,
% to find the structure of the basis 
% associated with~$T_3$ we apply
% the maps $t-3$ and~$(t-3)^2$, the maps $t-2$ and~$(t-2)^2$, and 
% the map~$t+1$.
% For instance,
% the basis has a total of three elements that are mapped to~$\zero$
% when we apply $t-3$ twice.
% It has two elements
% that are sent to~$\zero$ on two applications of~$t-2$.
% And it has one element sent to~$\zero$ by~$t+1$.
% That is, to find the elements of this basis we want to
% iterate the maps, so we look in  
% the generalized null spaces $\gennullspace{t-3}$, $\gennullspace{t-2}$,
% and~$\gennullspace{t+1}$.

接下来的例子都从一个矩阵出发，
算出与之相似的一个 Jordan 形式矩阵。
在此之前，我们再说明一点。
例如，在下面的第一个例子中我们将求出特征值为 $2$
与~$6$。
为挑选与~\( 2\) 相关的基元素，我们要求出
$t-2$、\( (t-2)^2 \) 与~\( (t-2)^3\) 的零空间。
也就是说，为找出这个例子的基中的元素，我们要在广义零空间
$\gennullspace{t-2}$ 与~$\gennullspace{t-6}$ 中寻找。
但我们必须确保对每个映射的计算不会影响对另一个的计算。

% So before those examples, 
% suppose that $\map{t}{\C^n}{\C^n}$ is a transformation and 
% that $\lambda\in\C$ is an eigenvalue.
% Note that the generalized null space~$\gennullspace{t-\lambda}$ 
% and range space~$\genrangespace{t-\lambda}$ are $t-\lambda$
% invariant spaces.
% For the generalized null space, 
% if $\vec{v}\in\gennullspace{t-\lambda}$ then 
% $(t-\lambda)^n(\vec{v})=\zero$ and it follows that
% $(t-\lambda)^n(\,(t-\lambda)(\vec{v})\,)=(t-\lambda)^{n+1}(\vec{v})=\zero$ 
% also. 
% As to the generalized range space,
% Lemma Five.III.\ref{le:RangeAndNullChains} and its proof show that
% $\genrangespace{t-\lambda}
%   =\rangespace{(t-\lambda)^n}
%   =\rangespace{(t-\lambda)^{n+1}}$ 
% and so $\vec{v}\in\genrangespace{t-\lambda}$ implies that
% $(t-\lambda)(\vec{v})\in\genrangespace{t-\lambda}$ also.

% Before the theorem, we first will show that applying $t-\lambda$ 
% won't do something awkward to the basis elements not in 
% $\gennullspace{t-\lambda}$.
% For instance, in the prior example we want that applying
% $t-3$ won't do something awkward to~$\vec{\beta}_4$, 
% $\vec{\beta}_5$, and~$\vec{\beta}_6$. 
% Here, awkward means that
% where \( \map{t}{V}{V} \) is a linear transformation,
% the restriction\appendrefs{restrictions of functions}\spacefactor=1000  %
% of \( t \) to a subspace \( M\subseteq V \) 
% need not be a linear transformation on \( M \):~there
% may be an \( \vec{m}\in M \)
% with \( t(\vec{m})\not\in M \). 
% An example is where $t$ is the transformation 
% that rotates the plane by
% a quarter turn and $M$ is the $x=y$ line; the vector $\binom{1}{1}$
% is rotated away from the subspace.
% We next show this does not happen in the special case of 
% eigenvalues;~if $j\neq i$ then 
% application of \( t-\lambda_j \) leaves 
% \( t-\lambda_i \)'s part unchanged.

\begin{definition} \label{def:invariant}
设 \( \map{t}{V}{V} \) 是一个变换。
若子空间 \( M\subseteq V \) 满足：只要 \( \vec{m}\in M \) 就有
\( t(\vec{m})\in M \)（简记为 \( t(M)\subseteq M \)），
则称它是 \definend{$t$ 不变的}%
\index{invariant subspace}\index{subspace!invariant}
\end{definition}

\begin{lemma} \label{le:tInvIfftMinLambdaInv}
子空间是 \( t \) 不变的当且仅当
它对任意标量 \( \lambda \) 都是 \( t-\lambda \) 不变的。
特别地，
若 \( \lambda_i,\lambda_j\in\C \) 是
\( t \) 的两个不相等的特征值，则
空间 \( \gennullspace{t-\lambda_i} \)
与 \( \genrangespace{t-\lambda_i} \)
都是 \( t-\lambda_j \) 不变的。
\end{lemma}

\begin{proof}
第一句中从右到左的方向是平凡的：~若子空间对
所有标量 $\lambda$ 都是 $t-\lambda$ 不变的，
则取 $\lambda=0$
即知它是 $t$ 不变的。
对于另一个方向，设子空间是 $t$~不变的，
即若 $\vec{m}\in M$ 则 $t(\vec{m})\in M$，并设 $\lambda$ 为
一个标量。
子空间 $M$ 对线性组合封闭，因此若
$t(\vec{m})\in M$ 则 $t(\vec{m})-\lambda\vec{m}\in M$。
于是若 $\vec{m}\in M$ 则 $(t-\lambda)\,(\vec{m})\in M$。

第二句由第一句得到。
这两个空间是 $t-\lambda_i$~不变的，所以它们是 \( t \)~不变的。
再用一次第一句就得到
它们也是 \( t-\lambda_j \) 不变的。
\end{proof}

\begin{example} \label{ex:FirstJordForm}
矩阵
\begin{equation*}
   T=
   \begin{mat}[r]
     2  &0  &1  \\
     0  &6  &2  \\
     0  &0  &2
   \end{mat}
\end{equation*}
的特征多项式为 \( (x-2)^2(x-6) \)。
\begin{equation*}
   |T-xI|=
   \begin{vmat}
     2-x  &0    &1  \\
     0    &6-x  &2  \\
     0    &0    &2-x
   \end{vmat}
   =(x-2)^2(x-6)
\end{equation*}
我们先处理特征值~$2$。
计算 $T-2I$ 的幂以及相应的零空间与零度
都是常规计算。
（回忆我们的约定：$T$ 表示关于标准基的
映射 $\map{t}{\C^3}{\C^3}$。）
\begin{center}
  \renewcommand{\arraystretch}{1.25}
  \begin{tabular}{r|ccc} 
    \multicolumn{1}{c}{\( p \)}  
         &\( (T-2I)^p \) &\( \nullspace{(t-2)^p}  \) 
         &\textit{零度}                                            \\  
    \hline
    \( 1 \)
    &\matrixvenlarge{\begin{mat}[r]
          0  &0  &1  \\
          0  &4  &2  \\
          0  &0  &0
        \end{mat}}
    &\( \set{\matrixvenlarge{\colvec{x \\ 0 \\ 0}}
         \suchthat x\in\C}  \)  
    &$1$                                                   \\
    \( 2 \)
    &\matrixvenlarge{\begin{mat}[r]
          0  &0  &0  \\
          0  &16 &8  \\
          0  &0  &0
        \end{mat}}
    &\( \set{\matrixvenlarge{\colvec{x \\ -z/2 \\  z}}
               \suchthat x,z\in\C}  \) 
    &$2$                                                   \\
    \( 3 \)
    &\matrixvenlarge{\begin{mat}[r]
          0  &0  &0  \\
          0  &64 &32 \\
          0  &0  &0
        \end{mat}}
    &\textit{--同上--}
    &\textit{--同上--}
  \end{tabular}
\end{center}
所以广义零空间 $\gennullspace{t-2}$ 的维数为二。
我们知道 $t-2$ 在这个子空间上的限制是幂零的。
由零度的增长方式可知，$t-2$ 在某个链基上的作用是
$\vec{\beta}_1\mapsto\vec{\beta}_2\mapsto\zero$。
于是我们可以把这个限制表示成标准形
\begin{equation*}
  N_2=
  \begin{mat}[r]
    0  &0  \\
    1  &0  
  \end{mat}
  =\rep{t-2}{B,B}
  \qquad
   B_2=\sequence{\colvec[r]{1 \\ 1 \\ -2},
                 \colvec[r]{-2 \\ 0 \\ 0}}  
\end{equation*}
（也可以选取别的基）。
因此，$t$ 在 $\gennullspace{t-2}$ 上的限制的作用
就由这个 Jordan 块表示。
\begin{equation*}
  J_2=N_2+2I=\rep{t}{B_2,B_2}=
  \begin{mat}[r]
    2  &0  \\
    1  &2
  \end{mat}
\end{equation*}

第二个特征值是 $6$。
对它的计算比较容易。
因为 $x-6$ 在特征多项式中的幂是一次的，
所以 $t-6$ 在 $\gennullspace{t-6}$ 上的限制
必是幂零的，幂零指数为一
（它不可能有小于一的指数，又因为 $x-6$ 是特征多项式中
指数为一的因子，它也不可能有大于一的指数）。
它在链基上的作用必是 $\vec{\beta}_3\mapsto\zero$，并且
由于它是零映射，其标准形 $N_6$
是 $\nbyn{1}$ 零矩阵。
因此，$t$ 在 $\gennullspace{t-6}$ 上作用的
标准形 $J_6$ 是单个元素为 $6$ 的
$\nbyn{1}$ Jordan 块
矩阵。
至于基，我们可以取广义零空间中任意一个非零向量。
\begin{equation*}
   B_6=\sequence{\colvec[r]{0 \\ 1 \\ 0}}
\end{equation*}

把这两部分合在一起就得到，
\( T \) 的 Jordan 形式是
\begin{equation*}
   \rep{t}{B,B}=
   \begin{mat}[r]
      2  &0  &0  \\
      1  &2  &0  \\
      0  &0  &6
   \end{mat}
\end{equation*}
其中 \( B \) 是 $B_2$ 与 $B_6$ 的连接。
（如果我们想谨慎地得到唯一的 Jordan 形式，那么
比如我们可以重排基元素，让特征值
按降序排列。
不过，升序排列也可以。）
\end{example}

\begin{example}  \label{SecJordanForm}
这个矩阵
与前一个例子有相同的特征多项式
\( (x-2)^2(x-6) \)。
\begin{equation*}
   T=
   \begin{mat}[r]
     2  &2  &1  \\
     0  &6  &2  \\
     0  &0  &2
   \end{mat}
\end{equation*}
但这里
$t-2$ 的作用在仅应用一次之后就稳定了\Dash 即
$t-2$ 在 $\gennullspace{t-2}$ 上的限制是幂零指数为一的幂零映射。
\begin{center}
  % \renewcommand{\arraystretch}{1.25}
  \begin{tabular}{r|ccc}
    \multicolumn{1}{c}{\( p \)}  &\( (T-2I)^p \)  &\( \nullspace{(t-2)^p}  \)  
       &\textit{零度} \\                                     \hline
   \( 1 \)
   &\matrixvenlarge{\begin{mat}[r]
         0  &2  &1  \\
         0  &4  &2  \\
         0  &0  &0
       \end{mat}}
   &\( \set{\matrixvenlarge{\colvec{x \\ (-1/2)z \\ z}}\suchthat x,z\in\C}  \) 
   &$2$                                                     \\
   \( 2 \)
   &\matrixvenlarge{\begin{mat}[r]
         0  &8  &4 \\
         0  &16 &8 \\
         0  &0  &0
       \end{mat}}
   &\textit{--同上--}
   &\textit{--同上--}
 \end{tabular}
\end{center}
所以 $t-2$ 在广义零空间上的限制在某个链基上
通过两条链起作用：$\vec{\beta}_1\mapsto\zero$ 
以及 $\vec{\beta}_2\mapsto\zero$。
我们有这个与特征值~$2$ 相关联的 Jordan 块。
\begin{equation*}
  J_2=
  \begin{mat}[r]
    2  &0  \\
    0  &2  
  \end{mat}
\end{equation*}

这样，与前一个例子的对比在于：虽然
特征多项式告诉我们要考察
$t-2$ 在其广义零空间上的作用，
但特征多项式并没有完全描述 $t-2$ 的作用。
我们必须做一些计算才能发现
极小多项式是 \( (x-2)(x-6) \)。

对于特征值 $6$，前一个例子中对第二个特征值的论证
在这里同样适用：~因为 $x-6$ 在
特征多项式中的幂是一次的，
所以 $t-6$ 在 $\gennullspace{t-6}$ 上的限制
必是幂零指数为一的幂零映射。
或者，我们也可以直接算出来。
\begin{center}
  \begin{tabular}{r|ccc}
    \multicolumn{1}{c}{\( p \)}  &\( (T-6I)^p \)  &\( \nullspace{(t-6)^p}  \)  
       &\textit{零度} \\                                     \hline
   \( 1 \)
   &\matrixvenlarge{\begin{mat}[r]
        -4  &2  &1  \\
         0  &0  &2  \\
         0  &0  &-4
       \end{mat}}
   &\( \set{\matrixvenlarge{\colvec{x \\ 2x \\ 0}}\suchthat x\in\C}  \) 
   &$1$                                                     \\
   \( 2 \)
   &\matrixvenlarge{\begin{mat}[r]
        16  &-8 &-4 \\
         0  &0  &-8 \\
         0  &0  &16
       \end{mat}}
   &\textit{--同上--}
   &\textit{--同上--}
 \end{tabular}
\end{center}
无论哪种方式，在 $\gennullspace{t-6}$ 上，限制 $t-6$ 的
标准形 $N_6$ 都是 $\nbyn{1}$ 零矩阵。
Jordan 块 $J_6$ 是元素为 $6$ 的 $\nbyn{1}$ 矩阵。
 
因此，$T$ 的 Jordan 形式是一个对角矩阵，所以~\( t \) 是一个
可对角化的变换。
\begin{equation*}
  \rep{t}{B,B}=
  \begin{mat}[r]
    2  &0  &0  \\
    0  &2  &0  \\
    0  &0  &6
  \end{mat}
  \qquad
  B=\cat{B_2}{B_6}
   =\sequence{\colvec[r]{1 \\ 0 \\ 0},
              \colvec[r]{0 \\ 1 \\ -2},
              \colvec[r]{1 \\ 2 \\ 0}}
\end{equation*}
在这三个基向量中，前两个来自
$t-2$ 的零空间，第三个来自~$t-6$ 的零空间。
\end{example}

\begin{example} \label{ThirdJordanForm}
对
\begin{equation*}
   T=
   \begin{mat}[r]
     -1  &4  &0  &0  &0  \\
      0  &3  &0  &0  &0  \\
      0  &-4 &-1 &0  &0  \\
      3  &-9 &-4 &2  &-1 \\
      1  &5  &4  &1  &4
   \end{mat}
\end{equation*}
做一些计算可知，
其特征多项式
是 \( (x-3)^3(x+1)^2 \)。
这张表
\begin{center}
  % \renewcommand{\arraystretch}{1.25}
  \begin{tabular}{@{}r|c@{}c@{}c@{}}
    \multicolumn{1}{c}{\( p \)}  
            &\( (T-3I)^p \)  &\( \nullspace{(t-3)^p}  \)
             &\textit{零度}   
    \\  \hline
    \( 1 \)
    &\matrixvenlarge{\begin{mat}[r]
         -4  &4  &0  &0  &0  \\
          0  &0  &0  &0  &0  \\
          0  &-4 &-4 &0  &0  \\
          3  &-9 &-4 &-1 &-1 \\
          1  &5  &4  &1  &1
       \end{mat}}
    &\( \set{\matrixvenlarge{\colvec{-(u+v)/2 \\
                     -(u+v)/2 \\
                      (u+v)/2 \\
                       u      \\
                       v}}\suchthat u,v\in\C}  \) 
    &$2$                                            \\
    \( 2 \)
    &\matrixvenlarge{\begin{mat}[r]
         16  &-16&0  &0  &0  \\
          0  &0  &0  &0  &0  \\
          0  &16 &16 &0  &0  \\
        -16  &32 &16 &0  &0  \\
          0  &-16&-16&0  &0
       \end{mat}}
    &\( \set{\matrixvenlarge{\colvec{ -z      \\
                      -z      \\
                       z      \\
                       u      \\
                       v}}\suchthat z,u,v\in\C}  \) 
    &$3$                                              \\
    \( 3 \)
    &\matrixvenlarge{\begin{mat}[r]
        -64  &64   &0   &0   &0  \\
          0  &0    &0   &0   &0  \\
          0  &-64  &-64 &0   &0  \\
         64  &-128 &-64 &0   &0  \\
          0  &64   &64  &0   &0
       \end{mat}}
    &\textit{--同上--}
    &\textit{--同上--}
  \end{tabular}
\end{center}
这表明 $t-3$ 在 $\gennullspace{t-3}$ 上的限制在某个
链基上通过两条链起作用：
$\vec{\beta}_1\mapsto\vec{\beta}_2\mapsto\zero$
以及
$\vec{\beta}_3\mapsto\zero$。

对另一个特征值做类似的计算
\begin{center}
  \renewcommand{\arraystretch}{1.25}
  \begin{tabular}{r|ccc}
    \multicolumn{1}{c}{\( p \)}  
         &\( (T+1I)^p \)  &\( \nullspace{(t+1)^p}  \) 
         &\textit{零度}  \\  
    \hline
    \( 1 \)
    &\matrixvenlarge{\begin{mat}[r]
          0  &4  &0  &0  &0  \\
          0  &4  &0  &0  &0  \\
          0  &-4 &0  &0  &0  \\
          3  &-9 &-4 &3  &-1 \\
          1  &5  &4  &1  &5
       \end{mat}}
    &\( \set{\matrixvenlarge{\colvec{-(u+v)   \\
                       0      \\
                      -v      \\
                       u      \\
                       v}}\suchthat u,v\in\C}  \)  
    &$2$                                              \\
    \( 2 \)
    &\matrixvenlarge{\begin{mat}[r]
          0  &16 &0  &0  &0  \\
          0  &16 &0  &0  &0  \\
          0  &-16&0  &0  &0  \\
          8  &-40&-16&8  &-8 \\
          8  &24 &16 &8  &24
       \end{mat}}
    &\textit{--同上--}
    &\textit{--同上--}
  \end{tabular}
\end{center}
得到 $t+1$ 在其广义零空间上的限制
在某个链基上通过两条单独的链起作用：
$\vec{\beta}_4\mapsto\zero$ 与 $\vec{\beta}_5\mapsto\zero$。

因此
\( T \) 相似于这个 Jordan 形式矩阵。
\begin{equation*}
   \begin{mat}[r]
     -1  &0  &0  &0  &0  \\
      0  &-1 &0  &0  &0  \\    
      0  &0  &3  &0  &0  \\
      0  &0  &1  &3  &0  \\
      0  &0  &0  &0  &3  
   \end{mat}
\end{equation*}
\end{example}





